\chapter{From Procedure to Description}\label{chap:1}

\begin{quote}
\emph{``Civilization advances by extending the number of important
operations which we can perform without thinking about them.''} ---
Alfred North Whitehead
\end{quote}

A software engineer's instinct on a combinatorial problem is to design
an algorithm: pick a strategy --- greedy, backtracking, branch-and-bound,
dynamic programming --- and commit to the search structure, the data
structures, the order of exploration, the pruning rules. The skill is
real and worth keeping. Many problems are best solved by writing the
procedure that solves them.

On most real combinatorial problems, however, a hand-rolled procedure is
far from the state of the art. Industrial solvers represent decades of
accumulated engineering --- clause learning, constraint propagation,
restart strategies, parallel workers, instance-tuned heuristics --- that
no individual programmer can reproduce in a project. The pragmatic move
on a hard combinatorial problem is therefore not to write a solver, but
to \emph{describe} the problem in a form that an existing solver can
consume, and let that solver search.

This is a paradigm shift, not a swap of one technique for another. The
two mindsets answer different questions:

\begin{itemize}
\item
  The \textbf{procedural} view answers \emph{how do I compute an
  answer?} It is the perspective of an algorithm designer. The output is
  a procedure that, when run, produces a solution.
\item
  The \textbf{declarative} view answers \emph{what counts as an answer?}
  It is the perspective of a problem describer. The output is a
  \emph{model} --- a precise, executable statement of which
  configurations are acceptable and which is best.
\end{itemize}

Both views are legitimate, and a working optimization practitioner moves
between them. Most readers of this manuscript will have stronger code
intuition than mathematical-notation intuition; that intuition does not
become useless under the new paradigm --- most of the work is still code
--- but it has to be redirected, since the code you write here
\emph{describes} rather than \emph{executes}.

The solver we will commit to in this manuscript is \textbf{CP-SAT}, the
constraint-programming solver that ships inside Google's
\href{https://developers.google.com/optimization}{OR-Tools}. The choice
is not innocent: different solvers reward different idioms, and a model
written for CP-SAT will not be byte-for-byte the model you would write
for a mixed-integer programming solver or a SAT solver. The underlying
paradigm shift, and most of the modeling vocabulary, transfers without
change. The solver-specific details are local to the last few chapters.

\section{A Concrete Contrast: 0/1 Knapsack}\label{sec:1.1}

The classical \textbf{0/1 knapsack} makes the shift concrete. Given a
set of items \(I\), each with a weight \(w_i\) and a value \(c_i\), and
a knapsack of capacity \(C\), choose a subset of items whose total
weight does not exceed \(C\) and whose total value is as large as
possible. The problem is a useful contrast because both an algorithmic
recipe and a declarative formulation are short enough to fit on one
page.

\subsection{The procedural answer}\label{sec:1.1.1}

The textbook procedural solution is a dynamic program. We define a
table \lstinline!dp[i][w]! whose entry is the best total value
achievable using the first \(i\) items and capacity at most \(w\), and
we fill the table by deciding, for each \((i, w)\), whether to take
item \(i\) or leave it.

\begin{lstlisting}[language=Python]
def knapsack(weights, values, C):
    n = len(weights)
    dp = [[0] * (C + 1) for _ in range(n + 1)]
    for i in range(1, n + 1):
        for w in range(C + 1):
            dp[i][w] = dp[i - 1][w]
            if weights[i - 1] <= w:
                take = dp[i - 1][w - weights[i - 1]] + values[i - 1]
                if take > dp[i][w]:
                    dp[i][w] = take
    return dp[n][C]
\end{lstlisting}

Several non-trivial decisions sit inside that function. We chose the
recursion (take versus leave), the table layout (items by capacity),
the fill order (forward in \(i\), forward in \(w\)), and the off-by-one
indexing between the table and the input arrays. None of these
decisions is part of the \emph{problem}; all of them are part of one
possible \emph{solution}.

The running time is proportional to \(\lvert I \rvert \cdot C\), which
is \emph{pseudo-polynomial}: linear in the numerical value of \(C\),
not in the bits needed to write \(C\) down. When \(C\) reaches the
millions the table no longer fits in memory, and the algorithm becomes
unusable even for trivially small \(\lvert I \rvert\). Other
algorithmic approaches exist, but they tend to be even more complex to
implement.

\subsection{The declarative answer}\label{sec:1.1.2}

The same problem written as mathematics is three lines:

\[
\begin{aligned}
\max \quad      & \sum_{i \in I} c_i\, x_i \\
\text{s.t.} \quad   & \sum_{i \in I} w_i\, x_i \;\le\; C \\
                    & x_i \in \{0, 1\} && \forall i \in I.
\end{aligned}
\]

Read aloud: maximize the total value over a choice of items, subject to
the total weight not exceeding the capacity, where each \(x_i\) is a
binary decision (take or leave). Three lines describe the entire problem
and commit to no algorithm at all. This is the form a paper or a
textbook would use, and \Cref{chap:2} develops the notation needed to
read it fluently.

The translation from this math into a working CP-SAT model is almost
line-for-line:

\begin{lstlisting}[language=Python]
from ortools.sat.python import cp_model

model = cp_model.CpModel()
x = [model.new_bool_var(f"x_{i}") for i in I]

model.add(sum(w[i] * x[i] for i in I) <= C)
model.maximize(sum(c[i] * x[i] for i in I))

solver = cp_model.CpSolver()
solver.solve(model)
\end{lstlisting}

One Boolean variable per item, one linear inequality for the capacity,
one linear objective, one call to the solver. There is no loop over the
capacity, no table, no decision about the order of exploration. CP-SAT
does all of that internally. The gain, for this example, is primarily
\emph{modeling style}: the code mirrors the mathematical statement
directly and is easy to extend. For knapsack specifically, specialized
algorithms remain very strong baselines, so the point of the example is
the declarative form, not a claim that CP-SAT is the canonical fastest
method for this problem family.

\subsection{Comparing the two}\label{sec:1.1.3}

The two formulations differ in three ways that matter.

The first is performance. The dynamic program's \(\lvert I \rvert
\cdot C\) running time is fine for moderate capacities and bad for
large ones, while
CP-SAT inherits decades of constraint-solver research --- clause
learning, propagation, restarts, parallel workers --- none of which we
had to write. This does not mean CP-SAT is automatically the best method
for knapsack; on narrow problem classes a tailored algorithm is often
the right baseline. The declarative model gives access to the solver's
machinery without requiring its implementation.

The second is reading effort. The procedural version mixes the problem
with its solution: to read it six months from now, you have to
reconstruct why the recursion is what it is, why the loops run in this
order, and what each index means. The declarative version reads as a
problem statement --- every line corresponds to a fact about the
problem, and the model is reviewable and explainable in roughly the same
form a colleague would expect on a whiteboard.

The third is extensibility, and it is the deepest of the three. Suppose
a new requirement appears: items A and B are mutually exclusive, or at
least three of items C, D, E must be taken. In the dynamic program the
recursion changes, the table changes, possibly the time complexity
changes, and we are back at the whiteboard. In the declarative model
each new requirement is one extra line --- \(x_A + x_B \le 1\) for the
mutual exclusion, \(x_C + x_D + x_E \ge 3\) for the cardinality --- and
the solver absorbs the addition without any change to how the search
proceeds. Real problems mutate; their constraints arrive in waves, first
the one in the textbook, then the regulatory rule, then the customer's
preference, then the operational quirk. A model that absorbs new
constraints by adding lines is incomparably easier to maintain than a
procedure that has to be redesigned each time.

\section{Modeling, Not Magic}\label{sec:1.2}

Handing a problem to a solver does not make it easy. NP-hard problems
remain NP-hard; the solver does not abolish complexity, it merely
searches more cleverly than you would have. What changes with
declarative modeling is not whether the problem is hard but how large an
\emph{instance} you can practically solve.

This last point is where the real skill of optimization modeling lives.
The same problem can be encoded in many mathematically equivalent ways,
and the encodings differ enormously in how well a solver handles them. A
naive encoding may time out at fifty variables; an idiomatic encoding
may solve in two seconds at five thousand. The two models are saying the
same thing about the same problem, but only one of them is saying it in
a way the solver can exploit.

A few patterns recur. Solvers are tuned for certain shapes of
constraint, so a formulation that reaches for the solver's native
vocabulary --- global constraints, indicator variables, well-known
reformulations --- usually beats one that reinvents the same logic from
arithmetic primitives. Encodings that rule out symmetric or equivalent
solutions, removing duplicates the solver would otherwise explore,
outperform encodings that leave that freedom in, even when both are
mathematically correct. None of this is visible from the problem
statement; choosing the representation is part of the modeler's job, and
two correct models for the same problem can differ in solving time by
orders of magnitude.

The remaining chapters cover the basics to get you rolling with building
and solving your first models. Performance engineering of models is a
topic for later.

\section{What This Manuscript Develops}\label{sec:1.3}

Four ideas run through the rest of the book, and each one is the subject
of one or more chapters.

The first is that \emph{mathematical notation is a tool, not
gatekeeping}. Optimization papers and textbooks state problems in
notation because notation is shorter and less ambiguous than prose ---
not because the community is exclusive. Reading and writing it is a
habit anyone working with optimization can pick up. \Cref{chap:2} is a fast
alignment layer for that notation: sets, variables, expressions,
predicates, and the bridges between the numeric and Boolean parts of the
language.

The second is that \emph{optimization models decompose into five parts}:
the entities the problem talks about, the parameters that describe a
particular instance, the decision variables, the constraints that
restrict their combinations, and the objective that ranks the feasible
ones. This decomposition is the structure every model in the manuscript
uses. \Cref{chap:3} walks through it in a from-scratch recipe.

The third is that \emph{a verifier is an orthogonal contract on what the
problem actually means}. A verifier is a small piece of plain Python
that takes a candidate solution and answers two questions: is it valid,
and how good is it? It does not produce solutions; it judges them.
Writing one is independent of which solver eventually does the search,
and it pins the problem definition down in code so later modeling
decisions cannot drift away from it. \Cref{chap:4} develops the verifier
pattern; later chapters treat it as optional but useful.

The fourth is that \emph{knowing the solver's language matters}. The
same problem can be modeled in many mathematically equivalent ways, and
an idiomatic encoding for one solver looks different from an idiomatic
encoding for another. This manuscript commits to \textbf{CP-SAT}, the
constraint-programming solver that ships inside Google's
\href{https://developers.google.com/optimization}{OR-Tools}. CP-SAT has
an expressive vocabulary --- global constraints for routing, scheduling,
packing, and combinatorial structure --- and \Cref{chap:5}--\Cref{chap:8} lay it out: a
first end-to-end model, then variables, constraints, and objectives in
reference form. \Cref{chap:9} closes the manuscript by reflecting on these
four ideas and naming a few practical moves that come up in real
modeling work.

\begin{platypusinfo}
This manuscript works directly with CP-SAT's Python API
rather than with an algebraic modeling language such as AMPL, GAMS,
Pyomo, MiniZinc, or CPMpy. Algebraic modeling languages add a layer of
abstraction above solver APIs and often allow the same model to target
multiple solvers. They are valuable tools, but they would obscure the
solver-specific idioms this manuscript wants to teach. CP-SAT's Python
API is already high-level enough for introductory modeling, while still
exposing the modeling choices that matter for performance.
\end{platypusinfo}
